\begin{table}[H]
\centering
\caption{Standardized Effect Sizes for Main Outcomes}
\label{tab:sde}
\small
\begin{threeparttable}
\begin{tabular}{lcccccl}
\toprule
Outcome & $\hat{\beta}$ & SE & SD($Y$) & SDE & SE(SDE) & Classification \\
\midrule
\multicolumn{7}{l}{\textit{Panel A: Pooled}} \\
Overall IMR & 0.024 & 0.471 & 10.60 & 0.0023 & 0.0444 & Null \\
Amenable MR & -0.269 & 0.313 & 7.43 & -0.0362 & 0.0421 & Small negative \\
Non-amenable MR & 0.101 & 0.257 & 3.23 & 0.0312 & 0.0793 & Small positive \\
Neonatal MR & 0.123 & 0.380 & 6.44 & 0.0191 & 0.0590 & Small positive \\
\midrule
\multicolumn{7}{l}{\textit{Panel B: Heterogeneous (baseline IMR split)}} \\
Amenable MR (high baseline) & -0.231 & 0.461 & 7.61 & -0.0304 & 0.0606 & Small negative \\
Amenable MR (low baseline) & -0.176 & 0.680 & 4.45 & -0.0395 & 0.1530 & Small negative \\
\bottomrule
\end{tabular}
\begin{tablenotes}[flushleft]
\scriptsize
\item[] \parbox{\linewidth}{\tolerance=9999\textit{Notes:} \textbf{Country:} Mexico. \textbf{Research question:} Does Seguro Popular, Mexico's non-contributory public health insurance program rolled out to municipalities in staggered waves from 2002 to 2005, reduce cause-specific infant mortality in covered municipalities? \textbf{Policy mechanism:} Seguro Popular provides free access to a package of essential health services for the uninsured population, including prenatal care, institutional delivery, neonatal intensive care, and treatment of childhood diarrheal and respiratory diseases, thereby reducing financial barriers to healthcare utilization for pregnant women and infants. \textbf{Outcome definition:} Cause-specific infant mortality rates (deaths of children under age 1 per 1,000 estimated live births), classified by ICD-10 cause of death into amenable causes (perinatal conditions P00--P96, diarrheal disease A00--A09, respiratory infections J00--J22, vaccine-preventable A33--A37) and non-amenable causes (congenital malformations Q00--Q99, external causes V01--Y98). \textbf{Treatment:} Binary indicator for state-level Seguro Popular enrollment, with four staggered cohorts (2002, 2003, 2004, 2005). \textbf{Data:} INEGI/Secretar\'{i}a de Salud death microdata (individual death records with ICD-10 cause codes), 1998--2012, at the municipality-year level; 1,404 municipalities across 15 years ($N$ = 20,998). \textbf{Method:} Callaway--Sant'Anna (2021) doubly robust staggered difference-in-differences with not-yet-treated controls and standard errors clustered at the state level (32 clusters). \textbf{Sample:} Mexican municipalities with at least 50 mean annual deaths (excluding very small municipalities with extremely noisy mortality rates); all 32 states included. SDE $= \hat{\beta} / \text{SD}(Y)$ where SD($Y$) is the unconditional standard deviation of the outcome variable. Classification refers to magnitude, not statistical significance: Large ($|$SDE$| > 0.15$), Moderate ($0.05$--$0.15$), Small ($0.005$--$0.05$), Null ($< 0.005$).}
\end{tablenotes}
\end{threeparttable}
\end{table}
